% ---------------- 第4章 章末 ----------------
\zhangmo{本章小结}
\subsection*{一、知识总结}
\begin{center}
\small
\begin{tabularx}{\linewidth}{p{2.0cm}p{3.6cm}X}
\toprule
平衡 & 常数 & 要点 \\
\midrule
电离平衡 & $K_\text{a}=\frac{c(\ce{H+})c(\ce{A-})}{c(\ce{HA})}$ & 稀释一般促进电离；升温影响看热效应；同条件下 $K_\text{a}$ 大酸性强；强酸制弱酸 \\
水的电离 & $K_\text{w}=c(\ce{H+})c(\ce{OH-})$ & 酸碱抑制、水解促进、升温促进；酸碱性看 $c(\ce{H+})$ 与 $c(\ce{OH-})$ 的相对大小 \\
盐类水解 & $K_\text{h}=K_\text{w}/K_\text{a}$ & 越弱越水解；酸式盐比较 $K_{\text{a}2}$ 与 $K_{\text{h}2}$；水解应用 \\
沉淀溶解 & $K_\text{sp}=c^m(\ce{A^{n+}})c^n(\ce{B^{m-}})$ & $Q$ 与 $K_\text{sp}$ 比较；沉淀完全的标准按题设，常取 $10^{-5}$ mol·L$^{-1}$；转化常数按方程式组合 \\
\bottomrule
\end{tabularx}
\end{center}

\subsection*{二、方法提升}
\begin{enumerate}
\item 遇到比值变化题，把比值改写成“$K$/某浓度”或“$K_1/K_2$”，再看那个浓度怎么变。
\item 做粒子浓度题，先写出三大守恒；比较大小时先找主要过程（电离还是水解），再定次序；等式题优先用守恒组合检验。
\item 做图像题，找交点、特殊点代入 $K$ 表达式；对数图注意斜率的意义。
\item 做 $K_\text{sp}$ 计算，开始沉淀用原浓度，沉淀完全用 $10^{-5}$；比较不同离子谁先沉淀要分别计算所需 $c(\ce{OH-})$（或 pH）。
\end{enumerate}

\subsection*{三、易错点归纳}
\begin{enumerate}
\item 把导电性强弱等同于电解质强弱；把难溶等同于弱电解质。
\item 认为稀释弱酸时所有离子浓度都减小（$c(\ce{OH-})$ 增大）。
\item 电荷守恒漏乘电荷数；物料守恒漏写某种型体。
\item 用 pH$=7$ 判断中性而忽略温度；把“水电离的 $c(\ce{H+})$”与溶液的 $c(\ce{H+})$ 混淆。
\item 不同类型难溶物直接用 $K_\text{sp}$ 比较溶解度；认为 $K_\text{sp}$ 小的物质不能转化为 $K_\text{sp}$ 大的物质。
\item 用总水解式代替多元弱酸根的分步水解式，掩盖第一步占主导的事实；高中书写通常要求分步。
\end{enumerate}

\zhangmo{分层练习}
\zu{A组\quad 基础巩固}
\begin{xiti}[start=1]
\item 常温下，将 0.1 mol·L$^{-1}$ \ce{CH3COOH} 溶液加水稀释，下列各量增大的是\xz
\xx{$c(\ce{H+})$}{$c(\ce{CH3COOH})$}{$\dfrac{c(\ce{CH3COO-})}{c(\ce{CH3COOH})}$}{$K_\text{a}$}

\item 下列说法正确的是\xz
\xxa{25 ℃ 时，某溶液中由水电离出的 $c(\ce{H+})=10^{-12}$ mol·L$^{-1}$，该溶液一定呈酸性}{100 ℃ 时纯水的 pH$\approx6.1$，此时纯水呈酸性}{向水中加入少量 \ce{NH4Cl} 固体，水的电离平衡正向移动，$K_\text{w}$ 不变}{向水中加入少量 NaOH 固体，水的电离平衡逆向移动，$c(\ce{H+})$ 增大}

\item 常温下，有 pH 均为 3 的盐酸和醋酸各 10 mL，下列说法正确的是\xz
\xxa{两溶液的物质的量浓度相等}{分别与足量锌反应，醋酸产生的 \ce{H2} 多}{分别加水稀释 100 倍，盐酸的 pH 比醋酸小}{分别与等浓度 NaOH 溶液反应至中性，消耗 NaOH 一样多}

\item 下列事实与盐类水解\textbf{无关}的是\xz
\xxa{\ce{NH4Cl} 溶液可用于除铁锈}{实验室盛放 \ce{Na2CO3} 溶液的试剂瓶不能用磨口玻璃塞}{用 \ce{Na2S} 溶液除去废水中的 \ce{Hg^2+}}{加热蒸干 \ce{AlCl3} 溶液得不到 \ce{AlCl3} 固体}

\item 对 0.1 mol·L$^{-1}$ \ce{Na2S} 溶液，下列关系正确的是\xz
\xxa{$c(\ce{Na+})=2c(\ce{S^2-})$}{$c(\ce{Na+})+c(\ce{H+})=c(\ce{S^2-})+c(\ce{HS-})+c(\ce{OH-})$}{$c(\ce{OH-})=c(\ce{H+})+c(\ce{HS-})+2c(\ce{H2S})$}{$c(\ce{Na+})>c(\ce{S^2-})>c(\ce{HS-})>c(\ce{OH-})$}

\item 已知 25 ℃ 时 $K_\text{sp}(\ce{AgCl})=1.8\times10^{-10}$。下列说法正确的是\xz
\xxa{AgCl 在水中和在 0.1 mol·L$^{-1}$ NaCl 溶液中的溶解度相同}{向 AgCl 饱和溶液中加入少量 NaCl 固体，$K_\text{sp}(\ce{AgCl})$ 减小}{将 $2\times10^{-4}$ mol·L$^{-1}$ \ce{AgNO3} 溶液与 $2\times10^{-4}$ mol·L$^{-1}$ NaCl 溶液等体积混合，有沉淀生成}{AgCl 饱和溶液中 $c(\ce{Ag+})=1.8\times10^{-10}$ mol·L$^{-1}$}

\item 按高中要求分步书写水解方程式，下列正确的是\xz
\xxa{\ce{CO3^2- + 2H2O <=> H2CO3 + 2OH-}}{\ce{Fe^3+ + 3H2O = Fe(OH)3 v + 3H+}}{\ce{NH4+ + H2O <=> NH3*H2O + H+}}{\ce{HCO3- + H2O <=> CO3^2- + H3O+}（表示 \ce{HCO3-} 的水解）}

\item 25 ℃ 时：(1) 0.050 mol·L$^{-1}$ \ce{Ba(OH)2} 溶液的 pH$=$\tk；(2) pH$=2$ 的盐酸稀释到原体积的 1000 倍，pH$=$\tk；(3) pH$=3$ 的盐酸与 pH$=11$ 的 NaOH 溶液等体积混合，pH$=$\tk；(4) 0.010 mol·L$^{-1}$ 一元弱酸 HA 溶液中 $c(\ce{H+})=1.0\times10^{-4}$ mol·L$^{-1}$，HA 的 $K_\text{a}\approx$\tk。
\end{xiti}

\zu{B组\quad 高考典型}
\begin{xiti}
\item \laiyuan{原创}\ 25 ℃ 时：\ce{H2CO3} $K_{\text{a}1}=4.5\times10^{-7}$、$K_{\text{a}2}=4.7\times10^{-11}$；HClO $K_\text{a}=4.0\times10^{-8}$；HCN $K_\text{a}=6.2\times10^{-10}$；\ce{CH3COOH} $K_\text{a}=1.75\times10^{-5}$。(1) 等浓度的 \ce{CH3COONa}、NaClO、NaCN、\ce{Na2CO3} 溶液，pH 由大到小的顺序为\tk。(2) 写出向 NaCN 溶液中通入少量 \ce{CO2} 的离子方程式。(3) 0.1 mol·L$^{-1}$ NaClO 溶液的 pH 约为\tk（$\sqrt{2.5}\approx1.6$，$\lg1.6\approx0.2$）。

\item 用 0.1000 mol·L$^{-1}$ NaOH 标准溶液滴定 25.00 mL 未知浓度的醋酸。(1) 应选用的指示剂是\tk，终点现象为\tk。(2) 三次滴定消耗 NaOH 溶液体积分别为 20.02 mL、19.98 mL、20.80 mL，醋酸的浓度为\tk。(3) 下列操作使测定结果偏高的是\tk（填序号）：①碱式滴定管未用标准液润洗；②锥形瓶用蒸馏水洗后未干燥；③滴定前仰视读数、滴定后平视读数；④滴定前滴定管尖嘴有气泡，滴定后气泡消失。

\item \laiyuan{改编自2024年广东卷第19题(3)}\ 一定温度下，某研究组用分光光度法测定两种一元弱酸 HX（X 为 A 或 B）在某非水溶剂中的 $K_\text{a}$。a．选择合适的指示剂 HIn，$K_\text{a}(\ce{HIn})=3.6\times10^{-20}$；其钾盐为 KIn。b．向 KIn 溶液中加入 HX，发生反应 \ce{In- + HX <=> X- + HIn}。KIn 起始物质的量为 $n_0(\ce{KIn})$，加入 HX 的物质的量为 $n(\ce{HX})$，平衡时测得 $c_\text{平}(\ce{In-})/c_\text{平}(\ce{HIn})$ 随 $n(\ce{HX})/n_0(\ce{KIn})$ 的变化：当 $n(\ce{HX})/n_0(\ce{KIn})=1.0$ 时，HA 体系中该比值为 3.0，HB 体系中该比值小于 1.0。已知该溶剂本身不电离，钾盐在该溶剂中完全电离；忽略游离 \ce{H+} 对物料衡算的影响。①计算 $K_\text{a}(\ce{HA})$（写出计算过程，结果保留两位有效数字）。②在该溶剂中，$K_\text{a}(\ce{HB})$\tk$K_\text{a}(\ce{HA})$；$K_\text{a}(\ce{HB})$\tk$K_\text{a}(\ce{HIn})$（填“$>$”“$<$”或“$=$”）。

\item \laiyuan{改编自2026年广东卷第17题(4)}\ 探究铜氨配合物 $Q$（\ce{[Cu(NH3)4]SO4*H2O}）与酸的反应。资料：向铜氨溶液中滴加硫酸，依次发生 (a)\ce{[Cu(NH3)4]^2+ + 2H+ + 2H2O = Cu(OH)2 v + 4NH4+}；(b)\ce{Cu(OH)2 + 2H+ = Cu^2+ + 2H2O}。实验：(ii) 称取 $w$ g Q 于烧杯中，移取 $V_0$ mL 稀氨水使 Q 完全溶解，得待测液；(iii) 搅拌下向待测液中滴加 0.1 mol·L$^{-1}$ \ce{H2SO4} 溶液，记录 pH–$V$ 数据；(iv) 移取 $V_0$ mL 上述稀氨水，用\tk{}溶液滴至终点，记录体积为 $V_1$ mL。①(iii) 所得 pH–$V$ 曲线呈现两次 pH 突跃，对应终点体积依次为 $V_2$ mL 和 $V_3$ mL，则反应(a)、(b)消耗酸的体积比 $V_\text{a}:V_\text{b}=$\tk（用 $V_1$、$V_2$、$V_3$ 表示）。②多次实验均得 $V_\text{a}:V_\text{b}>1$。有人提出假设：(iii) 中生成的沉淀除 \ce{Cu(OH)2} 外还含有 \ce{Cu2(OH)2SO4}。某次实验 $V_\text{a}:V_\text{b}=21:19$，则第一个终点时 $\dfrac{n[\ce{Cu(OH)2}]}{n[\ce{Cu2(OH)2SO4}]}=$\tk。

\item \laiyuan{原创}\ 某废液中含 0.10 mol·L$^{-1}$ \ce{Zn^2+} 和少量 \ce{Fe^3+}、\ce{Fe^2+}。已知 $K_\text{sp}[\ce{Fe(OH)3}]=2.8\times10^{-39}$，$K_\text{sp}[\ce{Fe(OH)2}]=4.9\times10^{-17}$，$K_\text{sp}[\ce{Zn(OH)2}]=1.0\times10^{-17}$，$\lg2\approx0.3$，$\lg6.5\approx0.8$。(1) 为什么要先加 \ce{H2O2} 再调 pH？写出相关离子方程式。(2) 计算使 \ce{Fe^3+} 沉淀完全而 \ce{Zn^2+} 不沉淀的 pH 范围。(3) 可选用什么物质调节 pH？

\item \laiyuan{原创}\ 常温下，将 0.2 mol·L$^{-1}$ 氨水与 0.1 mol·L$^{-1}$ 盐酸等体积混合（$K_\text{b}(\ce{NH3*H2O})=1.8\times10^{-5}$）。(1) 混合后溶液中的溶质及其浓度；(2) 判断溶液的酸碱性并说明理由；(3) 将下列粒子浓度由大到小排序：\ce{NH4+}、\ce{Cl-}、\ce{NH3*H2O}、\ce{OH-}、\ce{H+}；(4) 写出该溶液的物料守恒式。

\item \laiyuan{原创}\ 已知 $K_\text{sp}(\ce{AgCl})=1.8\times10^{-10}$，$K_\text{sp}(\ce{AgI})=8.5\times10^{-17}$。(1) 向含 AgCl 沉淀的饱和溶液中加入 KI 溶液，发生的反应为\tk，其平衡常数 $K=$\tk。(2) 若要使 0.01 mol AgCl 完全转化为 AgI（溶液体积 1 L），所加 KI 的物质的量至少为多少（结果保留两位有效数字）？
\end{xiti}

\zu{C组\quad 冲刺高分}
\begin{xiti}
\item \laiyuan{原创}\ 常温下，向 \ce{H2A} 溶液中滴加 NaOH 溶液，\ce{H2A}、\ce{HA-}、\ce{A^2-} 的分布系数 $\delta$ 随 pH 变化：\ce{H2A} 与 \ce{HA-} 曲线交于 pH$=1.2$，\ce{HA-} 与 \ce{A^2-} 曲线交于 pH$=4.2$。(1) 求 $K_{\text{a}1}$、$K_{\text{a}2}$。(2) NaHA 溶液呈酸性还是碱性？通过计算说明。(3) 当 $c(\ce{H2A})=c(\ce{A^2-})$ 时，溶液 pH 为多少？(4) 求 $\dfrac{c^2(\ce{HA-})}{c(\ce{H2A})\cdot c(\ce{A^2-})}$ 的值，它是否随 pH 变化？(5) 溶液 pH$=7$ 时，写出 $c(\ce{Na+})$ 与含 A 粒子浓度的关系式。

\item \laiyuan{改编自2026年广东卷第18题(3)}\ 金精矿提金流程中，“沉锑”用铁粉还原 \ce{[SbCl4]-} 得到 Sb，“沉锑”后的滤液含 \ce{Fe^2+}、\ce{Mg^2+} 等。“电解”在用阴离子交换膜隔开的电解池中进行，一定条件下实现“沉锑”试剂（铁粉）的回收利用。忽略生成气体的反应及溶液体积变化，设电解终了阳极室的 \ce{Fe^2+} 全部转化为 \ce{Fe^3+}。已知 $K_\text{sp}$（298 K）：\ce{Fe(OH)3} $2.8\times10^{-39}$，\ce{Fe(OH)2} $4.9\times10^{-17}$，\ce{Mg(OH)2} $5.6\times10^{-12}$。①阴极反应式为\tk。②初始电解液中 $c(\ce{Fe^2+})=1.4$ mol·L$^{-1}$，$c(\ce{Mg^2+})=0.056$ mol·L$^{-1}$。为防止阳极室电解过程中产生沉淀，理论上最终阳极室 $c(\ce{OH-})$ 应不大于\tk mol·L$^{-1}$（列算式）。

\item \laiyuan{原创}\ 用 0.1000 mol·L$^{-1}$ 盐酸滴定 20.00 mL 0.1000 mol·L$^{-1}$ \ce{Na2CO3} 溶液，忽略 \ce{CO2} 逸出，pH 计记录的滴定曲线有两个突跃：第一个终点在 20.00 mL 附近（pH 约 8.3），第二个终点在 40.00 mL 附近（pH 约 3.9）；加入 10.00 mL 时 pH$=10.33$，加入 30.00 mL 时 pH$=6.35$。(1) 写出两个阶段的离子方程式。(2) 求 \ce{H2CO3} 的 $K_{\text{a}1}$、$K_{\text{a}2}$（用 10 的幂表示）。(3) 若用指示剂代替 pH 计，两个终点应分别选用什么指示剂？(4) 加入 20.00 mL 盐酸时，溶液中的主要溶质是什么？写出此时的质子守恒式。
\end{xiti}

\zhangmo{练习解析}
\begin{jiexilist}
\item \ans{C}稀释时 $c(\ce{H+})$、$c(\ce{CH3COOH})$ 减小；$\dfrac{c(\ce{CH3COO-})}{c(\ce{CH3COOH})}=\dfrac{K_\text{a}}{c(\ce{H+})}$，$c(\ce{H+})$ 减小，比值增大；$K_\text{a}$ 只与温度有关，不变。
\item \ans{C}A：水的电离被抑制，可能是酸也可能是碱。B：纯水中 $c(\ce{H+})=c(\ce{OH-})$，始终呈中性。C：\ce{NH4+} 水解促进水的电离，温度不变 $K_\text{w}$ 不变，正确。D：加 NaOH，$c(\ce{OH-})$ 增大，$c(\ce{H+})$ 减小。
\item \ans{B}pH 相同，醋酸浓度远大于盐酸，A 错；醋酸总物质的量大，产生 \ce{H2} 多，B 对；稀释 100 倍，盐酸 pH 变为 5，醋酸继续电离，pH 小于 5，C 错；醋酸消耗 NaOH 多，D 错。
\item \ans{C}A：\ce{NH4+} 水解呈酸性，可溶解铁锈；B：\ce{CO3^2-} 水解呈碱性，会与玻璃中 \ce{SiO2} 反应使瓶塞粘住；C：\ce{Hg^2+ + S^2- = HgS v}，是沉淀反应，与水解无关；D：\ce{AlCl3} 水解，HCl 挥发，蒸干得 \ce{Al(OH)3}。
\item \ans{C}A：\ce{S^2-} 水解，$c(\ce{Na+})>2c(\ce{S^2-})$，物料守恒应为 $c(\ce{Na+})=2[c(\ce{S^2-})+c(\ce{HS-})+c(\ce{H2S})]$。B：漏乘电荷数，应为 $2c(\ce{S^2-})$。C：质子守恒，正确。D：\ce{OH-} 来自两步水解和水的电离，$c(\ce{OH-})>c(\ce{HS-})$。
\item \ans{C}A：同离子效应，AgCl 在 NaCl 溶液中溶解度更小；B：$K_\text{sp}$ 只与温度有关，不变；C：混合后 $c(\ce{Ag+})=c(\ce{Cl-})=1\times10^{-4}$ mol·L$^{-1}$，$Q=1\times10^{-8}>K_\text{sp}=1.8\times10^{-10}$，有沉淀生成，正确；D：$c(\ce{Ag+})=\sqrt{1.8\times10^{-10}}\approx1.3\times10^{-5}$ mol·L$^{-1}$。
\item \ans{C}A：多元弱酸根应分步水解；B：水解程度小，用“\ce{<=>}”且不标沉淀；C 正确；D 表示的是 \ce{HCO3-} 的电离，水解应为 \ce{HCO3- + H2O <=> H2CO3 + OH-}。
\item (1) 13。$c(\ce{OH-})=2\times0.050=0.10$ mol·L$^{-1}$，$c(\ce{H+})=10^{-13}$ mol·L$^{-1}$。(2) 5。强酸稀释 $10^3$ 倍，$c(\ce{H+})$ 由 $10^{-2}$ 变为 $10^{-5}$ mol·L$^{-1}$，此时水电离的 \ce{H+} 仍可忽略。(3) 7。$n(\ce{H+})=10^{-3}V=n(\ce{OH-})$，恰好中和，NaCl 溶液呈中性。(4) $1.0\times10^{-6}$。$K_\text{a}=\dfrac{(10^{-4})^2}{0.010-10^{-4}}\approx1.0\times10^{-6}$（电离度 1\%，分母近似为 0.010）。
\item (1) \ce{Na2CO3}$>$NaCN$>$NaClO$>$\ce{CH3COONa}。对应酸的强弱：\ce{HCO3-}（$4.7\times10^{-11}$）$<$HCN（$6.2\times10^{-10}$）$<$HClO（$4.0\times10^{-8}$）$<$\ce{CH3COOH}（$1.75\times10^{-5}$），越弱越水解，碱性越强。（\ce{NaHCO3} 既水解又电离，其溶液 pH 约为 8.3 且几乎不随浓度变化，不能只凭 $K_{\text{a}1}$ 与其他盐排序，故本题不列入。）(2) \ce{CN- + CO2 + H2O = HCN + HCO3-}（HCN 的 $K_\text{a}$ 大于 \ce{HCO3-}，只能生成 \ce{HCO3-}）。(3) $K_\text{h}=\dfrac{10^{-14}}{4.0\times10^{-8}}=2.5\times10^{-7}$，$c(\ce{OH-})\approx\sqrt{2.5\times10^{-7}\times0.1}=\sqrt{2.5}\times10^{-4}\approx1.6\times10^{-4}$ mol·L$^{-1}$，pOH$\approx3.8$，pH$\approx10.2$。
\item (1) 酚酞；当滴入最后半滴 NaOH 溶液时，溶液由无色变为浅红色，且半分钟内不褪色。(2) 第三次数据（20.80 mL）与其余两次相差过大，应检查原因并补做；本题用前两次相近数据的平均值 20.00 mL 计算；$c=\dfrac{0.1000\times20.00}{25.00}=0.08000$ mol·L$^{-1}$。(3) ①④。①标准液被稀释，消耗体积偏大，偏高；②无影响；③滴定前仰视读数偏大，$V_\text{标}$ 测得值偏小，偏低；④气泡体积被计入，偏高。
\item ①设 KIn 与 HA 起始浓度均为 $c_0$，转化 $x$：平衡时 $c(\ce{In-})=c(\ce{HA})=c_0-x$，$c(\ce{A-})=c(\ce{HIn})=x$。由 $\dfrac{c_0-x}{x}=3.0$ 得 $x=0.25c_0$。该反应 $K_1=\dfrac{c(\ce{A-})c(\ce{HIn})}{c(\ce{In-})c(\ce{HA})}=\dfrac{c(\ce{A-})c(\ce{H+})}{c(\ce{HA})}\cdot\dfrac{c(\ce{HIn})}{c(\ce{In-})c(\ce{H+})}=\dfrac{K_\text{a}(\ce{HA})}{K_\text{a}(\ce{HIn})}$，而 $K_1=\dfrac{0.25c_0\times0.25c_0}{0.75c_0\times0.75c_0}=\dfrac19$，故 $K_\text{a}(\ce{HA})=\dfrac{3.6\times10^{-20}}{9}=4.0\times10^{-21}$。②$>$；$>$。HB 体系中 $\dfrac{c_0-y}{y}<1.0$，$y>0.5c_0$，$K_2=\dfrac{y^2}{(c_0-y)^2}>1>K_1$，所以 $K_\text{a}(\ce{HB})>K_\text{a}(\ce{HA})$；又 $K_2=\dfrac{K_\text{a}(\ce{HB})}{K_\text{a}(\ce{HIn})}>1$，$K_\text{a}(\ce{HB})>K_\text{a}(\ce{HIn})$。
\item (iv) 0.1 mol·L$^{-1}$ \ce{H2SO4}（与(iii)中滴定所用溶液相同，作空白对照，测出 $V_0$ mL 稀氨水自身消耗的酸）。①$(V_2-V_1):(V_3-V_2)$。第一个突跃对应反应(a)完成，此前消耗的酸包括稀氨水消耗的 $V_1$ 和反应(a)消耗的部分，$V_\text{a}=V_2-V_1$；第二个突跃对应反应(b)完成，$V_\text{b}=V_3-V_2$。②18。设 \ce{Cu(OH)2} $x$ mol、\ce{Cu2(OH)2SO4} $y$ mol：生成沉淀时 \ce{[Cu(NH3)4]^2+ + 2H+ + 2H2O = Cu(OH)2 v + 4NH4+}，\ce{2[Cu(NH3)4]^2+ + 6H+ + SO4^2- + 2H2O = Cu2(OH)2SO4 v + 8NH4+}，消耗 \ce{H+} $(2x+6y)$ mol；溶解沉淀时 \ce{Cu(OH)2 + 2H+ = Cu^2+ + 2H2O}，\ce{Cu2(OH)2SO4 + 2H+ = 2Cu^2+ + SO4^2- + 2H2O}，消耗 $(2x+2y)$ mol。$\dfrac{2x+6y}{2x+2y}=\dfrac{21}{19}$，解得 $x=18y$。
\item (1) 要把 \ce{Fe^2+} 降到 $10^{-5}$ mol·L$^{-1}$，需 $c(\ce{OH-})=\sqrt{4.9\times10^{-17}/10^{-5}}\approx2.2\times10^{-6}$ mol·L$^{-1}$，此时 \ce{Zn^2+} 会大量沉淀，不能这样除铁；先用 \ce{H2O2} 将 \ce{Fe^2+} 氧化为 \ce{Fe^3+}，\ce{Fe^3+} 在很低 pH 下即可沉淀完全：\ce{2Fe^2+ + H2O2 + 2H+ = 2Fe^3+ + 2H2O}。(2) \ce{Fe^3+} 沉淀完全：$c(\ce{OH-})=\sqrt[3]{2.8\times10^{-39}/10^{-5}}\approx6.5\times10^{-12}$，pOH$\approx11.2$，pH$\approx2.8$；\ce{Zn^2+} 开始沉淀：$c(\ce{OH-})=\sqrt{1.0\times10^{-17}/0.10}=10^{-8}$，pH$=6$。范围：$2.8\le$pH$<6$。(3) ZnO、\ce{Zn(OH)2} 或 \ce{ZnCO3}（消耗 \ce{H+} 且不引入新杂质）。
\item (1) 按完全中和先得到 \ce{NH4Cl} 和剩余 \ce{NH3*H2O}，两者的分析浓度均为 0.05 mol·L$^{-1}$；建立平衡后各粒子浓度会略有调整。(2) 碱性：\ce{NH3*H2O} 的电离常数 $K_\text{b}=1.8\times10^{-5}$ 大于 \ce{NH4+} 的水解常数 $K_\text{h}=\frac{10^{-14}}{1.8\times10^{-5}}\approx5.6\times10^{-10}$，电离程度大于水解程度。(3) $c(\ce{NH4+})>c(\ce{Cl-})>c(\ce{NH3*H2O})>c(\ce{OH-})>c(\ce{H+})$。(4) $c(\ce{NH4+})+c(\ce{NH3*H2O})=2c(\ce{Cl-})$。
\item (1) \ce{AgCl(s) + I-(aq) <=> AgI(s) + Cl-(aq)}；$K=\dfrac{c(\ce{Cl-})}{c(\ce{I-})}=\dfrac{1.8\times10^{-10}}{8.5\times10^{-17}}\approx2.1\times10^6$。(2) 完全转化后溶液中 $c(\ce{Cl-})=0.01$ mol·L$^{-1}$，平衡时 $c(\ce{I-})=\dfrac{0.01}{2.1\times10^6}\approx4.7\times10^{-9}$ mol·L$^{-1}$，可忽略；消耗 \ce{I-} 0.01 mol，所以至少加 KI 约 0.010 mol（理论上略多于 0.01 mol）。
\item (1) 交点处两种型体浓度相等：$K_{\text{a}1}=c(\ce{H+})=10^{-1.2}$，$K_{\text{a}2}=10^{-4.2}$。(2) 酸性。\ce{HA-} 的水解常数 $K_{\text{h}2}=\dfrac{10^{-14}}{10^{-1.2}}=10^{-12.8}<K_{\text{a}2}=10^{-4.2}$，电离程度大于水解程度。(3) $K_{\text{a}1}K_{\text{a}2}=\dfrac{c^2(\ce{H+})c(\ce{A^2-})}{c(\ce{H2A})}$，$c(\ce{H2A})=c(\ce{A^2-})$ 时 $c^2(\ce{H+})=10^{-5.4}$，pH$=2.7$。(4) $\dfrac{c^2(\ce{HA-})}{c(\ce{H2A})c(\ce{A^2-})}=\dfrac{K_{\text{a}1}}{K_{\text{a}2}}=10^{3}$，是常数，不随 pH 变化。(5) pH$=7$ 时 $c(\ce{H+})=c(\ce{OH-})$，由电荷守恒 $c(\ce{Na+})=c(\ce{HA-})+2c(\ce{A^2-})$。
\item ①\ce{Fe^2+ + 2e- = Fe}（回收铁粉，供“沉锑”再用）。②$\sqrt[3]{\dfrac{2.8\times10^{-39}}{1.4}}$。阳极上 \ce{Fe^2+} 失电子生成 \ce{Fe^3+}，最终阳极室 $c(\ce{Fe^3+})=1.4$ mol·L$^{-1}$。比较各离子开始沉淀所需 $c(\ce{OH-})$：\ce{Fe(OH)3} 为 $\sqrt[3]{2.8\times10^{-39}/1.4}\approx1.3\times10^{-13}$；\ce{Fe(OH)2}（若有剩余 \ce{Fe^2+}）约 $5.9\times10^{-9}$；\ce{Mg(OH)2} 为 $\sqrt{5.6\times10^{-12}/0.056}=10^{-5}$。\ce{Fe(OH)3} 所需最小，故 $c(\ce{OH-})\le\sqrt[3]{2.8\times10^{-39}/1.4}$ mol·L$^{-1}$。
\item (1) 第一阶段：\ce{CO3^2- + H+ = HCO3-}；第二阶段：\ce{HCO3- + H+ = H2CO3}（实际敞口滴定还可能有 \ce{CO2} 逸出，本题按不逸出处理）。(2) 加入 10.00 mL 时 \ce{CO3^2-} 恰好转化一半，$c(\ce{CO3^2-})\approx c(\ce{HCO3-})$，$K_{\text{a}2}\approx c(\ce{H+})=10^{-10.33}$；加入 30.00 mL 时 $c(\ce{HCO3-})\approx c(\ce{H2CO3})$，$K_{\text{a}1}\approx10^{-6.35}$。(3) 第一终点（pH$\approx$8.3）用酚酞；第二终点（pH$\approx$3.9）用甲基橙。(4) \ce{NaHCO3}（和 NaCl）；质子守恒：$c(\ce{H+})+c(\ce{H2CO3})=c(\ce{OH-})+c(\ce{CO3^2-})$。
\end{jiexilist}
